%% manuscript.tex — The Lancet format
%% Compilar com: pdflatex → bibtex → pdflatex × 2
%% ===================================================================
\documentclass[review,number]{elsarticle}

\usepackage[utf8]{inputenc}
\usepackage[T1]{fontenc}
\usepackage{amsmath,amssymb}
\usepackage{graphicx}
\usepackage{booktabs}
\usepackage{multirow}
\usepackage{xcolor}
\usepackage{lineno}
\usepackage{hyperref}
\usepackage[margin=2.5cm]{geometry}

\biboptions{super,comma}   % Superscript Vancouver citations (Lancet style)

\linenumbers  % Obrigatório para submissão

%% Metadados do journal
\journal{The Lancet}

\begin{document}

\begin{frontmatter}

%% ── Título ──────────────────────────────────────────────────────────
\title{The compounding respiratory hospitalisation risk of cold spells
       and particulate matter in a subtropical metropolis:
       a time-series analysis}

%% ── Autores ─────────────────────────────────────────────────────────
\author[ufpr]{Felipe Baglioli}
\author[utfpr]{Lucas Rover}
\author[utfpr]{Leonardo Jos\'{e} Rossoni Quadros}
\author[utfpr]{Yara de Souza Tadano}
\author[ufpr]{Eduardo Tadeu Bacalhau}
\author[ufpr]{Ricardo Henrique Moreton Godoi\corref{cor1}}
\ead{godoi@ufpr.br}
\cortext[cor1]{Corresponding author}

\address[ufpr]{Environmental Engineering Department, Federal University of Paran\'{a} (UFPR), Curitiba, PR, Brazil}
\address[utfpr]{Federal University of Technology -- Paran\'{a} (UTFPR), Ponta Grossa, PR, Brazil}

%% ── Summary (Lancet ≤300 palavras) ──────────────────────────────────
\begin{abstract}

\textbf{Background}\quad
Climate--health policy in Latin America has intensified the frequency
of compound environmental extremes; however, the joint health impacts
of cold spells concurrent with fine particulate matter
(PM$_{2{\cdot}5}$) accumulation remain understudied in subtropical
cities.

\textbf{Methods}\quad
We analysed daily respiratory hospital admissions (DATASUS SIH/AIH,
ICD-10 J00--J99, 2022--2024; $N$=42\,716), temperature (INMET,
1961--2024), and air quality (PM$_{2{\cdot}5}$) in Curitiba, Brazil
($\sim$1${\cdot}$8~million inhabitants) using a methodological
triangulation design. K-means clustering identified environmental
regimes and their lagged associations with hospitalisations;
distributed lag non-linear models (DLNMs) with quasi-Poisson
regression quantified the exposure--response relationship and lag
structure; and gradient-boosted trees (XGBoost) with SHAP feature
importance provided predictive convergence evidence.

\textbf{Findings}\quad
Over 1096~study days, 42\,716~respiratory hospitalisations (51\% male;
38\% aged $\geq$65~years) and 3715~in-hospital deaths were recorded.
Clustering identified a ``Clean Cold'' regime whose delayed correlations
with hospitalisations peaked at lags of three to 10~days. The DLNM
confirmed this pattern: MMT 17${\cdot}$8\,°C; cold extremes at the
1st percentile of $T_\mathrm{min}$ (5${\cdot}$2\,°C) yielded
cumulative RR 1${\cdot}$42 (95\% CI 1${\cdot}$09--1${\cdot}$84)
over zero to 21~days --- corresponding to approximately 14 excess
admissions per extreme cold day; at the 5th percentile
(8${\cdot}$5\,°C), RR 1${\cdot}$22 (1${\cdot}$02--1${\cdot}$44).
Heat effects were not statistically significant
(P90 RR 1${\cdot}$06, 0${\cdot}$98--1${\cdot}$14).
Results were robust across all sensitivity analyses
($\hat{\phi}$=1${\cdot}$32, Ljung--Box $P$=0${\cdot}$14).
SHAP analysis independently ranked $T_\mathrm{min}$ as the dominant
environmental predictor, completing the three-method convergence.

\textbf{Interpretation}\quad
In this subtropical metropolis, cold extremes --- not heat --- exhibit
the strongest association with temperature-related respiratory
hospitalisations, with a delayed effect structure that provides a three
to 10-day window for preventive intervention. Current early-warning
systems focus almost exclusively on heat alerts; integrating cold-spell
and air quality forecasts could reduce a substantial and preventable
respiratory burden.

\textbf{Funding}\quad
NAPI -- Arauc\'{a}ria Foundation, Paran\'{a}, Brazil.

\end{abstract}

%% ── Keywords ────────────────────────────────────────────────────────
\begin{keyword}
Environmental Epidemiology \sep Air Pollutants \sep
Public Health Surveillance \sep Extreme Cold \sep
Explainable Artificial Intelligence \sep Delayed Effects
\end{keyword}

\end{frontmatter}

%% ═══════════════════════════════════════════════════════════════════
%% RESEARCH IN CONTEXT
%% ═══════════════════════════════════════════════════════════════════

\section*{Research in Context}

\textbf{Evidence before this study}

We searched PubMed and Scopus for studies published between
January~2010 and December~2025 using the terms ``temperature
extremes'', ``hospital admissions'', ``distributed lag'', and
``Southern Hemisphere'' or ``Latin America''. The landmark
384-location study by Gasparrini et~al.\ (2015) established U-shaped
temperature--mortality curves globally, yet included no subtropical
South American city. Subsequent multi-city analyses have reinforced
that cold-attributable mortality exceeds heat-attributable mortality
in most climates, but morbidity evidence from the subtropical Latin
America remains scarce. Furthermore, existing studies have rarely
combined exploratory, inferential, and predictive approaches to
assess the joint role of temperature and air pollution on respiratory
morbidity. No study has applied a methodological triangulation design
--- integrating unsupervised clustering, DLNM, and machine learning
--- to quantify the non-linear, delayed effects of cold extremes and
PM$_{2{\cdot}5}$ on respiratory hospital admissions in this region.

\textbf{Added value of this study}

To our knowledge, this is the first study to evaluate the compounding
role of temperature extremes and air pollution on respiratory
hospitalizations in a subtropical South American city. We formally
quantified a 42\% excess risk of respiratory admissions during extreme
cold (1st percentile of $T_\mathrm{min}$, 5${\cdot}$2\,°C;
RR=1${\cdot}$42 [1${\cdot}$09--1${\cdot}$84]), while heat effects
were non-significant. More importantly, we identified a highly
actionable 3--10~day delayed window between cold exposure and the
hospitalization peak. The methodological convergence of our
independent frameworks (inferential and predictive machine learning
models) ensures this lag is a robust physiological and systemic
signal, challenging the heat-centric paradigm in Latin American
public health and providing a concrete operational timeframe for
early interventions.

\textbf{Implications of all the available evidence}

Early-warning systems in subtropical Latin American cities should
integrate cold-spell forecasts and air quality monitoring alongside
heat alerts. The three to 10-day delay between cold exposure and the
hospitalisation peak provides a concrete window for preventive action,
allowing primary health-care teams to proactively monitor patients at
higher risk, and emergency care units to preemptively adjust medical
stocks. Exploratory analyses consistently identified PM$_{2{\cdot}5}$
as a structural co-factor associated with higher admissions,
suggesting that combined temperature--pollution alerts may be more
effective than temperature-only triggers. Given that current Brazilian
climate--health surveillance (VIGIAR-SUS) lacks a cold-specific
trigger, our findings, corroborated by three independent analytical
approaches, identify a remediable gap in public health preparedness.

%% ═══════════════════════════════════════════════════════════════════
%% INTRODUCTION
%% ═══════════════════════════════════════════════════════════════════
\section{Introduction}

Global evidence increasingly shows that cold-attributable mortality
exceeds heat-attributable mortality in most
climates,\cite{Gasparrini2015} yet climate--health policy remains
overwhelmingly focused on heat alerts and heatwave preparedness.
This asymmetry is particularly pronounced in subtropical cities of
Latin America, where winter temperatures can fall below
5\,°C but early-warning systems rarely include cold-spell
triggers.\cite{Bell2024,McMichael2006}

The pathophysiology of cold-related respiratory morbidity is well
established. Low ambient temperatures suppress mucociliary clearance,
impair alveolar macrophage function, and increase indoor crowding ---
conditions that favour respiratory pathogen
transmission.\cite{Linares2025} Cold also triggers systemic
vasoconstriction, elevates blood viscosity, and raises cardiovascular
strain, compounding the respiratory burden in vulnerable
populations.\cite{Ebi2021} Critically, unlike heat-related morbidity,
which manifests acutely within 0--3~days, cold-related effects
exhibit delayed lags of 3--14~days, complicating their detection in
routine surveillance and their attribution in ecological
studies.\cite{Cheng2024} Concurrently, PM$_{2{\cdot}5}$ inhalation
consistently induces systemic inflammation and oxidative stress,
leading to endothelial dysfunction and autonomic
imbalance.\cite{Brook2010} The synchronisation of these stressors
imposes an additive burden on the respiratory system, where
cold-induced mucosal vulnerability is compounded by
particulate-driven inflammation.

Despite these known mechanisms, formal inferential evidence on the
temperature--morbidity relationship remains scarce for subtropical
South America. The landmark 384-location study by Gasparrini
et~al.\cite{Gasparrini2015} did not include any subtropical Southern
Brazilian city. Multi-city studies in Brazil have focused on air
pollution and mortality,\cite{Requia2024} while temperature--morbidity
analyses have relied on correlation-based approaches that cannot
disentangle delayed effects from shared seasonality. As a result,
a fundamental question remains unanswered for this region: does cold
or heat impose the greater respiratory hospitalization burden when
combined with high air pollution episodes, and over what time horizon?

To rigorously address these interactions, this study adopts the
framework of Compound Climate Extremes. According to Zscheischler
et~al.,\cite{Zscheischler2020} these events occur when a combination
of multiple drivers and/or hazards contributes to societal or
environmental risk. Crucially, the co-occurrence of these drivers
creates a severe concurrent exposure scenario that heavily compounds
the clinical burden on healthcare systems --- a dynamic particularly
critical when meteorological stagnation amplifies pollutant
concentrations during thermal extremes.

We address this gap through a methodological triangulation design
applied to 42\,716~respiratory hospital admissions, daily temperature
records, and EPA-corrected PM$_{2{\cdot}5}$ in Curitiba, the largest
subtropical city in southern Brazil. Curitiba's wide thermal range
provides a natural experiment for contrasting cold and heat effects
within a single population. Our primary objective is to quantify the
non-linear, delayed exposure--response relationship between
temperature extremes and respiratory admissions, with a pre-specified
hypothesis that cold effects would dominate in this setting. The
triangulation proceeds in three stages: unsupervised clustering first
explores the multivariate data structure to identify environmental
regimes and their lagged associations with hospitalisations without
parametric assumptions; distributed lag non-linear models (DLNM) then
formally quantify the exposure--response relationship and lag structure
suggested by the clustering; and gradient-boosted tree models with
SHAP feature-importance analysis offer an independent, predictive line
of convergence.

Beyond a local case study, Curitiba serves as a critical proxy for
densely populated metropolitan areas in the subtropical and temperate
transition zones of Latin America (eg, the Southern Cone: Southern
Brazil, Northern Argentina, Uruguay, Chile), which face similar rapid
urbanisation and housing insulation deficits. Unlike high-latitude
cities designed for low temperatures, this widespread structural
deficit theoretically renders populations paradoxically more
vulnerable to even moderate cold spells. This creates a hypothesised
``hidden vulnerability'' scenario where the health impact of
temperature drops is disproportionately compounded by co-occurring
air pollution, a dynamic often overlooked in Northern
Hemisphere-centric
literature.\cite{Gasparrini2015,Cohen2017}

%% ═══════════════════════════════════════════════════════════════════
%% METHODS
%% ═══════════════════════════════════════════════════════════════════
\section{Methods}

\subsection{Study design and setting}

We conducted an ecological time-series study of daily environmental
exposures and respiratory hospital admissions in Curitiba, Paraná,
Brazil (population~$\sim$1${\cdot}$8~million) from
1~January~2022 to 31~December~2024 (1096~days). Curitiba is the
largest city in southern Brazil with a Köppen Cfb (oceanic subtropical)
climate, winter minima below 5\,°C, and summer maxima approaching
30\,°C --- making it representative of the broader subtropical belt
of southern South America.

\subsection{Data sources}

Individual-level hospital admission records were obtained from the
Brazilian Unified Health System (SIH/SUS--DATASUS) for 2022--2024
($N$=42\,716 respiratory admissions, ICD-10 J00--J99). The primary
outcome was the aggregate daily count of all-cause respiratory
admissions. Demographics: 51\% of patients were of male sex; 38\%
were aged 65~years or older.

Daily minimum ($T_\mathrm{min}$), mean ($T_\mathrm{med}$), and
maximum ($T_\mathrm{max}$) temperatures, relative humidity, and
atmospheric pressure were obtained from INMET automatic station~A807.
$T_\mathrm{min}$ served as the primary exposure in the DLNM; all six
meteorological and pollution variables were used in the clustering and
machine learning analyses.

PM$_{2{\cdot}5}$ concentrations were measured using a single PurpleAir
Flex-II sensor (Sensor ID~90875) at daily resolution. Raw readings
were corrected using the EPA nationwide correction
equation:\cite{Barkjohn2021}
%
\begin{equation}
  \mathrm{PM}_{2{\cdot}5}^{\mathrm{corr}} =
    0{\cdot}524 \times \mathrm{PM}_{2{\cdot}5}^{\mathrm{sensor}}
    - 0{\cdot}0862 \times \mathrm{RH} + 5{\cdot}75
  \label{eq:epa}
\end{equation}
%
This correction has been validated against Federal Equivalent Method
monitors\cite{Holder2020} and is widely adopted for PurpleAir
data.\cite{Magi2020,Barkjohn2022} An inter-channel quality control
protocol excluded days when inter-channel differences exceeded
5~µg/m$^3$ and 70\% relative difference simultaneously. After QC,
697~days (63${\cdot}$6\%) had valid EPA-corrected PM$_{2{\cdot}5}$.
Crucially, the 36${\cdot}$4\% missingness was not missing completely
at random (not MCAR), occurring more frequently on colder days due to
seasonal sensor downtime, which fundamentally limits the statistical
power of subsequent inferential models (details in Supplement~S.1).

\subsection{Definitions and missing data}

Extreme temperature days were defined using ETCCDI monthly percentiles
(TX90p for extremely hot, TN10p for extremely cold) over the 1961--2024
baseline. Heatwaves: $\geq$3~consecutive extreme hot days; cold spells:
$\geq$3~consecutive extreme cold days. Missing data rates:
$T_\mathrm{min}$~2${\cdot}$6\%, PM$_{2{\cdot}5}$~36${\cdot}$4\%,
hospitalisations~0${\cdot}$2\%. Primary analyses used complete cases;
sensitivity analyses with imputation produced similar results
(Supplement).

\subsection{Clustering analysis (structural-exploratory evidence)}

As the first analytical step, K-means
clustering\cite{MacQueen1967} was applied to six standardised variables
($T_\mathrm{max}$, $T_\mathrm{med}$, $T_\mathrm{min}$,
PM$_{2{\cdot}5}$, relative humidity, atmospheric pressure) on
complete-case days, with optimal $K$ selected by silhouette
score.\cite{Rousseeuw1987} Lagged cross-correlation functions (lags
0--28~days) were computed between cluster membership and daily
hospitalisations to identify temporal associations without
parametric constraints. PCA was performed on standardised
PM$_{2{\cdot}5}$, $T_\mathrm{min}$, and hospitalisations ($N$=683).
This analysis does not impose a distributed lag function, spline basis,
or risk-ratio framework, and thus constitutes an independent,
assumption-free exploration of the lag structure.

\subsection{Distributed lag non-linear model (inferential evidence)}
\label{sec:dlnm}

Guided by the lag patterns identified in the clustering analysis,
distributed lag non-linear models
(DLNM)\cite{Gasparrini2010,Gasparrini2011} were applied to formally
quantify the non-linear exposure--response relationship and its
delayed lag structure:
%
\begin{equation}
  \log\!\big(\mathrm{E}[Y_t]\big) = \alpha
    + \mathbf{s}\bigl(x_t, \mathrm{lag}; \boldsymbol{\eta}\bigr)
    + \sum_{d=1}^{6} \beta_d \,\mathrm{DOW}_{d,t}
    + \gamma \,\mathrm{Holiday}_t
    + \mathrm{ns}(\mathrm{time}_t,\, 7\,\mathrm{df/year})
  \label{eq:dlnm}
\end{equation}
%
where $Y_t$ is the daily count of respiratory admissions (quasi-Poisson),
$\mathbf{s}(\cdot)$ is the cross-basis function (natural cubic splines:
exposure df=4, lag df=5), DOW denotes day-of-week indicators, and
$\mathrm{ns}(\mathrm{time}, 21\,\mathrm{df})$ captures trend and
seasonality. Maximum lag was set at 21~days, consistent with the
standard adopted in multi-country temperature--mortality
analyses\cite{Gasparrini2015} and reflecting the known 3--14-day
latency of cold-related respiratory
effects.\cite{Linares2025,Cheng2024} The minimum morbidity temperature
(MMT) --- defined as the temperature at which the cumulative
exposure--response curve reaches its minimum risk, analogous to the
minimum mortality temperature of Gasparrini
et~al.\cite{Gasparrini2015} --- was identified by grid search.
Sensitivity analyses varied the maximum lag (14, 21, 28~days),
seasonal df (6, 7, 8~df/year), and excluded 2022 to assess residual
COVID-19 effects.

\subsection{Machine learning (predictive convergence evidence)}

As a third, methodologically independent line of evidence, machine
learning models (XGBoost,\cite{Chen2016} multilayer
perceptron,\cite{Goodfellow2016} extreme learning
machine\cite{Huang2006}) were trained to forecast daily
hospitalisations at prediction horizons of 0, 1, 2, 3, 4, 5, 6, and
7~days using a 70/30 train--test split. The 7-day limit reflects the
practical ceiling of operational weather forecast accuracy and differs
from the epidemiological lag window (0--21~days) used in the DLNM to
capture the full delayed attribution of temperature effects: the ML
horizon asks ``how far ahead can we predict?'', whereas the DLNM lag
asks ``how long does the biological effect persist?''
SHAP values\cite{Lundberg2017} quantified feature importance. Because
gradient-boosted trees use recursive partitioning rather than spline
bases and optimise predictive accuracy rather than estimating relative
risks, agreement between SHAP rankings and DLNM findings constitutes
methodological convergence.

\paragraph{Exploratory scenario and extreme-event simulations}
To assess future health impacts under climate change, the trained
XGBoost and MLP models were applied to downscaled CMIP6
SSP5-8${\cdot}$5 temperature projections (EC-Earth3-Veg-LR) for
Curitiba's grid cell. Monte Carlo simulations of extreme events were
conducted by superimposing synthetic stressor episodes onto the 2050
climate baseline: cold waves (temperature drops of 34--88\% below
monthly means for 3--8~consecutive days) and wildfire PM$_{2{\cdot}5}$
spikes (70--100~µg/m$^3$ for 7--14~days).

\paragraph{Ethics and reporting}
This study used exclusively aggregated, publicly available data from
DATASUS, INMET, and PurpleAir. Under Brazilian CNS~510/2016,
research using public aggregated data is exempt from ethics review.
Reporting follows the STROBE guidelines for observational
studies.\cite{VonElm2007}

%% ═══════════════════════════════════════════════════════════════════
%% RESULTS
%% ═══════════════════════════════════════════════════════════════════
\section{Results}

\subsection{Descriptive statistics}

Over 1096~study days, we recorded 42\,716~respiratory hospital
admissions and 3715~in-hospital deaths
(case fatality 8${\cdot}$7\%),
representing a mean of 32${\cdot}$3~admissions per day
(SD~11${\cdot}$7, range 1--73; Table~\ref{tab:descriptive}). The
study population was predominantly male (51\%) and elderly (38\%
aged $\geq$65~years). EPA-corrected PM$_{2{\cdot}5}$ was generally
low (median~4${\cdot}$12~µg/m$^3$), with only 4${\cdot}$7\% of days
exceeding the WHO 24-hour guideline (15~µg/m$^3$). Using ETCCDI
monthly percentiles, we identified 215~extremely hot days
(20${\cdot}$1\%) and 39~extremely cold days (3${\cdot}$7\%),
aggregating into 25~heatwave events but only 4~cold-spell events
(eTable~1). Despite accounting for fewer than one-fifth of
extreme-temperature days, cold events were the only ones associated
with significant excess morbidity (see §\ref{sec:dlnm_results}) ---
a low-frequency, high-impact pattern that underscores the episodic
and concentrated nature of cold exposure in this climate and
highlights a gap in surveillance systems designed primarily for
frequent heat events.

\begin{table}[ht]
\centering
\caption{Descriptive statistics for all study variables, Curitiba,
2022--2024 ($N$=1096~days).}
\label{tab:descriptive}
\begin{tabular}{lrrrrrrrr}
\toprule
Variable & $N$ & Mean & SD & Median & Q25 & Q75 & Min & Max \\
\midrule
$T_\mathrm{max}$ (°C) & 1070 & 24${\cdot}$53 & 4${\cdot}$75 & 24${\cdot}$80 & 21${\cdot}$50 & 28${\cdot}$10 & 10${\cdot}$20 & 35${\cdot}$40 \\
$T_\mathrm{med}$ (°C) & 1042 & 18${\cdot}$74 & 3${\cdot}$63 & 18${\cdot}$90 & 16${\cdot}$30 & 21${\cdot}$50 &  7${\cdot}$80 & 27${\cdot}$30 \\
$T_\mathrm{min}$ (°C) & 1068 & 14${\cdot}$75 & 3${\cdot}$75 & 15${\cdot}$20 & 12${\cdot}$00 & 17${\cdot}$70 &  1${\cdot}$30 & 23${\cdot}$40 \\
PM$_{2{\cdot}5}$ (µg/m$^3$) &  697 &  5${\cdot}$84 & 7${\cdot}$40 &  4${\cdot}$12 &  2${\cdot}$43 &  6${\cdot}$55 &  0${\cdot}$00 & 65${\cdot}$19 \\
Rel.\ humidity (\%)     &  859 & 79${\cdot}$98 &  9${\cdot}$52 & 80${\cdot}$90 & 74${\cdot}$15 & 86${\cdot}$55 & 43${\cdot}$00 & 98${\cdot}$40 \\
Hosp.\ (n/day)   & 1094 & 32${\cdot}$29 & 11${\cdot}$74 & 32${\cdot}$00 & 23${\cdot}$00 & 41${\cdot}$00 & 1${\cdot}$00 & 73${\cdot}$00 \\
\bottomrule
\end{tabular}
\end{table}

\subsection{Clustering: structural-exploratory evidence}

K-means clustering on six climate--pollution variables ($K$=3,
silhouette=0${\cdot}$357) identified three environmental regimes:
Clean Cold ($n$=248), Polluted Heat ($n$=12), and Clean Heat
($n$=298; eTables~2--3, eFigures~1--4). Crucially, the Clean Cold
cluster exhibited delayed positive correlations with hospitalisations
peaking at lags 3--10~days without imposing any distributional lag
function or spline basis --- a pattern subsequently confirmed by the
formal DLNM analysis (§\ref{sec:dlnm_results}).
A validation analysis at $K$=4 yielded stronger between-cluster
differences in daily hospitalisations (Kruskal--Wallis
$H$=16${\cdot}$69, $p<$0${\cdot}$001; eTable~3): the Polluted Heat
cluster (PM$_{2{\cdot}5}$=51${\cdot}$5~µg/m$^3$, $n$=12) exhibited
the highest mean hospitalisations (38${\cdot}$6$\pm$9${\cdot}$7 per
day), suggesting that elevated PM$_{2{\cdot}5}$ compounds the
hospitalisation burden across environmental regimes.

PCA corroborated this pattern: PC1 (44${\cdot}$5\% of variance)
loaded positively on PM$_{2{\cdot}5}$ ($+$0${\cdot}$564) and
hospitalisations ($+$0${\cdot}$588) and negatively on
$T_\mathrm{min}$ ($-$0${\cdot}$580; eTable~4), indicating that cold,
polluted days are structurally associated with higher admissions.
Extended lag analysis confirmed peak
$T_\mathrm{min}$--hospitalisation correlation at lag~9
($r$=$-$0${\cdot}$320, $p<$0${\cdot}$001; eTable~5).

\subsection{Distributed lag non-linear models}
\label{sec:dlnm_results}

\paragraph{Temperature exposure--response}
The clustering analysis revealed delayed cold--hospitalisation
correlations peaking at lags 3--10~days. To formally quantify this
association, the quasi-Poisson DLNM for $T_\mathrm{min}$ ($N$=1068,
$\hat{\phi}$=1${\cdot}$32) identified a minimum morbidity temperature
(MMT) of 17${\cdot}$8\,°C, closely aligned with the 75th percentile
of the local temperature distribution
(Figure~\ref{fig:dlnm_temp_overall}). Below this threshold, the
cumulative exposure--response curve rose steeply and monotonically:
at the 1st percentile of $T_\mathrm{min}$ (5${\cdot}$2\,°C),
RR=1${\cdot}$42 [1${\cdot}$09--1${\cdot}$84] --- equivalent to a
42\% excess in respiratory admissions; at the 5th percentile
(8${\cdot}$5\,°C), RR=1${\cdot}$22 [1${\cdot}$02--1${\cdot}$44]; at
the 10th percentile (9${\cdot}$8\,°C), RR=1${\cdot}$17
[1${\cdot}$00--1${\cdot}$37] (Table~\ref{tab:dlnm_temp_rr}). In
stark contrast, heat effects were non-significant across all upper
percentiles (P90 RR=1${\cdot}$06 [0${\cdot}$98--1${\cdot}$14];
P99 RR=1${\cdot}$29 [0${\cdot}$98--1${\cdot}$68]). The lag
structure at cold (P10) peaked at 3--10~days
(Figure~\ref{fig:dlnm_temp_lag}a), confirming the delayed pattern
first identified by the clustering cross-correlations and consistent
with the known latency of cold-induced respiratory pathology.

\paragraph{Sensitivity and diagnostics}
The cold effect was robust to all pre-specified sensitivity analyses
(Table~\ref{tab:sensitivity}). Excluding 2022 --- the year with the
greatest residual COVID-19 influence --- strengthened the cold effect
to RR=1${\cdot}$37 [1${\cdot}$11--1${\cdot}$70], suggesting that
pandemic-related disruptions attenuated the temperature--morbidity
signal in the full period. Model diagnostics showed no residual
autocorrelation (Durbin--Watson=2${\cdot}$00, Ljung--Box
$p$=0${\cdot}$14; eTable~7).

\paragraph{Exploratory PM$_{2{\cdot}5}$ analysis}
An exploratory DLNM for EPA-corrected PM$_{2{\cdot}5}$ ($N$=697,
lag 0--14~days) showed no significant cumulative associations: at the
90th percentile (10${\cdot}$2~µg/m$^3$), RR=0${\cdot}$91
[0${\cdot}$78--1${\cdot}$07] (eFigure~6). A joint
temperature--PM$_{2{\cdot}5}$ model ($N$=683) yielded
RERI=0${\cdot}$013 [$-$0${\cdot}$006, 0${\cdot}$033], indicating
no supra-additive interaction on the additive scale (eTable~6).
To strictly avoid selective reporting, we emphasize this null
inferential finding. These null results should be interpreted in the
context of low ambient concentrations
(median~4${\cdot}$1~µg/m$^3$) and 36${\cdot}$4\% MNAR sensor data
missingness, which severely limited statistical power.

\begin{table}[ht]
\centering
\caption{Cumulative relative risk (RR) at key percentiles of
$T_\mathrm{min}$, centred at MMT=17${\cdot}$8\,°C. Quasi-Poisson DLNM,
lag 0--21~days, $N$=1068.}
\label{tab:dlnm_temp_rr}
\begin{tabular}{lrrr}
\toprule
Percentile & $T_\mathrm{min}$ (°C) & RR & 95\% CI \\
\midrule
P1  &  5${\cdot}$2 & 1${\cdot}$418 & 1${\cdot}$091--1${\cdot}$842 \\
P5  &  8${\cdot}$5 & 1${\cdot}$215 & 1${\cdot}$023--1${\cdot}$442 \\
P10 &  9${\cdot}$8 & 1${\cdot}$167 & 0${\cdot}$998--1${\cdot}$365 \\
P25 & 12${\cdot}$0 & 1${\cdot}$133 & 0${\cdot}$979--1${\cdot}$312 \\
P75 & 17${\cdot}$7 & 1${\cdot}$001 & 0${\cdot}$996--1${\cdot}$005 \\
P90 & 19${\cdot}$4 & 1${\cdot}$055 & 0${\cdot}$978--1${\cdot}$138 \\
P95 & 20${\cdot}$0 & 1${\cdot}$102 & 0${\cdot}$979--1${\cdot}$240 \\
P99 & 21${\cdot}$4 & 1${\cdot}$285 & 0${\cdot}$983--1${\cdot}$679 \\
\bottomrule
\end{tabular}
\end{table}

\begin{table}[ht]
\centering
\caption{Sensitivity analyses for the temperature DLNM. Cumulative RR
at the 10th percentile of $T_\mathrm{min}$ (9${\cdot}$8\,°C).}
\label{tab:sensitivity}
\begin{tabular}{lrrrr}
\toprule
Scenario & RR & 95\% CI & $\hat{\phi}$ & $N$ \\
\midrule
Main (lag=21, 7\,df/yr) & 1${\cdot}$17 & 1${\cdot}$00--1${\cdot}$37 & 1${\cdot}$32 & 1068 \\
Lag=14                  & 1${\cdot}$05 & 0${\cdot}$92--1${\cdot}$18 & 1${\cdot}$31 & 1068 \\
Lag=28                  & 1${\cdot}$23 & 1${\cdot}$01--1${\cdot}$51 & 1${\cdot}$32 & 1068 \\
Seasonal 6\,df/yr       & 1${\cdot}$17 & 1${\cdot}$00--1${\cdot}$36 & 1${\cdot}$33 & 1068 \\
Seasonal 8\,df/yr       & 1${\cdot}$21 & 1${\cdot}$03--1${\cdot}$42 & 1${\cdot}$31 & 1068 \\
Exclude 2022            & 1${\cdot}$37 & 1${\cdot}$11--1${\cdot}$70 & 1${\cdot}$26 &  718 \\
\bottomrule
\end{tabular}
\end{table}

\begin{figure}[ht]
\centering
\includegraphics[width=0.85\textwidth]{figures/fig_dlnm_temp_overall.pdf}
\caption{Cumulative exposure--response curve for $T_\mathrm{min}$
and respiratory hospitalisations (DLNM, lag 0--21~days). The curve
rises steeply below the MMT (17${\cdot}$8\,°C, red dashed line),
indicating a cold-dominant pattern. Heat effects above the MMT are
non-significant. Shaded area: 95\% CI.}
\label{fig:dlnm_temp_overall}
\end{figure}

\begin{figure}[ht]
\centering
\includegraphics[width=\textwidth]{figures/fig_dlnm_temp_lag.pdf}
\caption{Lag-specific RR for (a) cold extreme (P10, 9${\cdot}$8\,°C)
and (b) heat extreme (P90, 19${\cdot}$4\,°C), centred at
MMT=17${\cdot}$8\,°C.}
\label{fig:dlnm_temp_lag}
\end{figure}

\subsection{Machine learning: predictive convergence evidence}

\paragraph{Model performance across prediction horizons}
Three machine learning models were evaluated across eight prediction
horizons (lag 0--7~days; eTable~12, eFigure~10). XGBoost achieved
the best overall performance at the 1-day horizon
(MAPE=18${\cdot}$0\%, RMSE=8${\cdot}$10; eTable~8), outperforming
MLP (MAPE=23${\cdot}$2\%, RMSE=8${\cdot}$80) and ELM
(MAPE=28${\cdot}$1\%, RMSE=10${\cdot}$43). The fact that 1-day-ahead
predictions outperformed same-day predictions is consistent with the
delayed effect structure identified by the DLNM.

\paragraph{SHAP feature importance}
SHAP feature-importance analysis (eFigure~8) ranked features by mean
absolute SHAP value: weekday (0${\cdot}$54), temperature
(0${\cdot}$23), PM$_{2{\cdot}5}$ (0${\cdot}$16), humidity
(0${\cdot}$05), and holiday (0${\cdot}$03). Among environmental
predictors, temperature was the dominant feature and PM$_{2{\cdot}5}$
ranked second, contributing approximately 70\% of the predictive
importance of temperature. The beeswarm plot confirmed that low
temperature values and high PM$_{2{\cdot}5}$ values independently
increased predicted hospitalisations. This is noteworthy because
gradient-boosted trees use recursive partitioning --- a fundamentally
different approach from the DLNM's cross-basis splines --- yet both
temperature and PM$_{2{\cdot}5}$ emerge as the two leading
environmental predictors, with directional effects consistent across
methods.

\paragraph{Exploratory scenario simulations}
Monte Carlo simulations of compound extreme events superimposed onto
the 2050 climate baseline revealed substantial additional
hospitalisation burdens (eFigures~12--13). Cold wave scenarios
produced weekly hospitalisation increases of up to 17\% (MLP) and
13\% (XGBoost); wildfire PM$_{2{\cdot}5}$ scenarios
(70--100~µg/m$^3$ for 7--14~days) generated surges of up to 55\%
(MLP) and 42\% (XGBoost).

%% ═══════════════════════════════════════════════════════════════════
%% DISCUSSION
%% ═══════════════════════════════════════════════════════════════════
\section{Discussion}

\subsection{Principal findings}

This study provides formal DLNM-based evidence that cold
extremes --- not heat --- exhibit the strongest association with
temperature-related respiratory hospitalisations in a subtropical
South American city, and that PM$_{2{\cdot}5}$ operates as a
structural co-factor identified consistently across three independent
analytical approaches. At the 1st percentile of $T_\mathrm{min}$
(5${\cdot}$2\,°C), cumulative RR reached 1${\cdot}$42
[1${\cdot}$09--1${\cdot}$84] over 0--21~days, corresponding to
approximately 14 excess respiratory admissions per extreme cold day.
Heat effects, by contrast, were non-significant across all upper
percentiles, although extremely hot days were predominant in the
studied area (20${\cdot}$1\% against only 3${\cdot}$7\% extremely
cold days).

\subsection{Cold-dominant burden: challenging the heat-centric paradigm}

A dominant framing in climate--health research equates rising
temperatures with rising health
risk.\cite{Bell2024,McMichael2006} Our findings complicate this
narrative for subtropical settings. In Curitiba, the asymmetry
between cold and heat effects was striking: the cold extreme (P1)
yielded a statistically significant 42\% excess, while the heat
extreme (P90) produced a non-significant 6\% excess. This asymmetry
is further underscored by the frequency of extreme events: extremely
hot days accounted for 20${\cdot}$1\% of the study period compared
with only 3${\cdot}$7\% for extremely cold days, yet cold was the
only exposure associated with significant excess morbidity --- a
low-frequency, high-impact pattern that current surveillance systems
are not designed to detect.

Our contribution goes beyond confirming this pattern. First, we extend
the evidence from mortality to respiratory \emph{morbidity}. Second,
we employ a methodological triangulation design that has not been
applied in previous temperature--health studies, providing convergent
evidence from three independent analytical frameworks. Third, we
identify PM$_{2{\cdot}5}$ as a structural co-factor associated with
hospitalisations across all exploratory and predictive approaches
(PCA, clustering, ML/SHAP), suggesting a combined cold--pollution
pathway that warrants investigation in multi-city designs.

\subsection{Pathophysiological plausibility and lag structure}

The 3--10-day lag peak is biologically coherent. Cold exposure
initiates a cascade that unfolds over days: immediate airway cooling
suppresses mucociliary clearance and reduces epithelial barrier
integrity (hours); subsequent impairment of alveolar macrophage
phagocytosis and secretory IgA production increases susceptibility
to viral and bacterial infection (1--3~days); clinical illness and
hospitalisation then follow after an incubation and progression
period (3--10~days).\cite{Linares2025,Ebi2021} This temporal sequence
--- unlike the rapid onset of heat-related presentations --- explains
why cold-attributable morbidity is systematically underdetected by
surveillance systems that operate on a 0--48-hour event window.

Additionally, cold weather increases indoor crowding and reduces
ventilation, amplifying respiratory pathogen transmission at the
population level.\cite{Cheng2024} In subtropical Brazilian cities,
this behavioural pathway is amplified by a built-environment
vulnerability largely absent in temperate-climate regions with
comparable winter minima: residential buildings typically lack thermal
insulation and central heating.\cite{Gasparrini2015,Cohen2017}

\subsection{Robustness and policy implications}

The cold effect was robust to all pre-specified sensitivity analyses.
Notably, excluding 2022 strengthened the cold effect (RR=1${\cdot}$37
[1${\cdot}$11--1${\cdot}$70]), suggesting that pandemic-related
changes in healthcare-seeking behaviour may have attenuated the
temperature--morbidity signal. The methodological triangulation
strengthens the finding beyond any single approach: the three methods
differ in their assumptions, optimisation criteria, and outputs, yet
converge on the same conclusion.

These findings have direct implications for climate--health
preparedness in subtropical South America.\cite{Bell2024} In Brazil,
VIGIAR-SUS currently lacks a cold-specific trigger for respiratory
early warnings. The 3--10-day delayed peak provides a concrete
operational timeframe: days 0--2, primary health-care teams initiate
active tele-monitoring of vulnerable patients; days 3--5, emergency
care units preemptively reinforce stocks and adjust staffing.
However, as a single-city study, these results require validation in
multi-city pooled analyses before informing national
policy.\cite{Requia2024}

\subsection{Limitations}

Several limitations warrant consideration. First, as an ecological
time-series study, we cannot infer individual-level causal
relationships. Second, the 3-year analysis period encompasses only
three seasonal cycles; however, the DLNM estimates derive from the
full continuous temperature distribution ($N$=1068~days), and the
convergence of three independent approaches provides additional
assurance. Third, confounding by respiratory viral circulation
(influenza, RSV) cannot be excluded; without adjusting for weekly
local viral surveillance data, we cannot strictly disentangle direct
thermodynamic stress from weather-facilitated viral transmission.
Fourth, we used total respiratory admissions without age or
cause-specific stratification. Fifth, PM$_{2{\cdot}5}$ data had
36${\cdot}$4\% missing values (not MCAR); the structural associations
identified by PCA, clustering, and SHAP suggest that the
PM$_{2{\cdot}5}$ pathway may become inferentially detectable in
settings with higher pollution levels. Sixth, our 2050 simulations
hold demographics constant and serve to quantify climate-hazard
pressure, not as literal epidemiological forecasts. Finally, as a
single-city study, results require replication; however, the
methodological triangulation partially compensates by demonstrating
internal convergence.

\subsection{Conclusion}

In Curitiba, cold extremes --- not heat --- exhibit the strongest
association with temperature-related respiratory hospitalisations,
with a 42\% excess risk at the coldest days (RR=1${\cdot}$42
[1${\cdot}$09--1${\cdot}$84]) and a delayed effect peaking at
3--10~days that creates a concrete window for preventive
intervention. The convergence of three independent analytical
approaches --- structural-exploratory (clustering), inferential
(DLNM), and predictive (ML/SHAP) --- strengthens this finding by
demonstrating that the cold-lag pattern emerges from the data
regardless of modelling assumptions. These findings challenge the
heat-centric framing of climate--health surveillance in subtropical
South America and identify a specific, remediable gap: the absence
of cold-spell triggers in early-warning systems. Integrating a
temperature-threshold alert at the identified MMT
(17${\cdot}$8\,°C) into existing surveillance frameworks could
reduce a preventable respiratory burden in one of South America's
largest subtropical populations.

%% ═══════════════════════════════════════════════════════════════════
%% DECLARAÇÕES (Lancet backmatter order)
%% ═══════════════════════════════════════════════════════════════════
\section*{Contributors}
FB: Conceptualisation, Data curation, Formal analysis, Investigation,
Methodology, Writing --- original draft.
LR: Methodology, Software, Validation, Formal analysis, Writing ---
review \& editing.
LJRQ: Data curation, Validation.
ETB: Data curation, Formal analysis, Methodology, Writing --- review.
YST: Conceptualisation, Investigation, Formal analysis, Writing ---
review \& editing.
RHMG: Conceptualisation, Writing --- original draft, Supervision.
All authors had full access to all the data in the study and had final
responsibility for the decision to submit for publication.

FB and LR directly accessed and verified the underlying data reported
in this manuscript.

\section*{Declaration of interests}
The authors declare no competing interests.

\section*{Data sharing}
All code and processed datasets are available at
\url{https://github.com/Roverlucas/cwb-temp-pollution-health}.
Raw hospitalisation data: DATASUS
(\url{https://datasus.saude.gov.br/}); meteorological data: INMET
station~A807; PM$_{2{\cdot}5}$: PurpleAir Sensor ID~90875.

\section*{Funding}
This work was supported by the New Arrangements for Research and
Innovation (NAPI) programme of the Arauc\'{a}ria Foundation, Paran\'{a},
Brazil.

\section*{Acknowledgements}
The authors thank the Brazilian Institute of Meteorology (INMET),
DATASUS, and PurpleAir Inc.\ for providing open data.

%% ═══════════════════════════════════════════════════════════════════
%% REFERÊNCIAS (Vancouver, numbered by first citation)
%% ═══════════════════════════════════════════════════════════════════
\bibliographystyle{elsarticle-num}

\begin{thebibliography}{27}

\bibitem{Gasparrini2015}
Gasparrini A, Guo Y, Hashizume M, et al.
Mortality risk attributable to high and low ambient temperature: a multicountry observational study.
\textit{Lancet} 2015; \textbf{386}: 369--75.

\bibitem{Bell2024}
Bell ML, Gasparrini A, Benjamin GC.
Climate change, extreme heat, and health.
\textit{N Engl J Med} 2024; \textbf{390}: 1793--801.

\bibitem{McMichael2006}
McMichael AJ, Woodruff RE, Hales S.
Climate change and human health: present and future risks.
\textit{Lancet} 2006; \textbf{367}: 859--69.

\bibitem{Linares2025}
Linares C, et al.
How air pollution and extreme temperatures affect emergency hospital admissions due to various respiratory causes in Spain, by age group.
\textit{Int J Hyg Environ Health} 2025; \textbf{266}: 114570.

\bibitem{Ebi2021}
Ebi KL, et al.
Hot weather and heat extremes: health risks.
\textit{Lancet} 2021; \textbf{398}: 698--708.

\bibitem{Cheng2024}
Cheng C, et al.
Effects of extreme temperature events on deaths and its interaction with air pollution.
\textit{Sci Total Environ} 2024; \textbf{915}: 170212.

\bibitem{Brook2010}
Brook RD, Rajagopalan S, Pope CA 3rd, et al.
Particulate matter air pollution and cardiovascular disease: an update to the scientific statement from the American Heart Association.
\textit{Circulation} 2010; \textbf{121}: 2331--78.

\bibitem{Requia2024}
Requia WJ, et al.
Short-term air pollution exposure and mortality in Brazil: investigating the susceptible population groups.
\textit{Environ Pollut} 2024; \textbf{340}: 122797.

\bibitem{Zscheischler2020}
Zscheischler J, Martius O, Westra S, et al.
A typology of compound weather and climate events.
\textit{Nat Rev Earth Environ} 2020; \textbf{1}: 333--47.

\bibitem{Cohen2017}
Cohen AJ, Brauer M, Burnett R, et al.
Estimates and 25-year trends of the global burden of disease attributable to ambient air pollution: an analysis of data from the Global Burden of Diseases Study 2015.
\textit{Lancet} 2017; \textbf{389}: 1907--18.

\bibitem{Barkjohn2021}
Barkjohn KK, Gantt B, Clements AL.
Development and application of a United States-wide correction for PM2.5 data collected with the PurpleAir sensor.
\textit{Atmos Meas Tech} 2021; \textbf{14}: 4617--37.

\bibitem{Holder2020}
Holder AL, Mebust AK, Maghran LA, et al.
Field evaluation of low-cost particulate matter sensors for measuring wildfire smoke.
\textit{Sensors} 2020; \textbf{20}: 4796.

\bibitem{Magi2020}
Magi BI, Cupini C, Francis J, et al.
Evaluation of PM2.5 measured in an urban setting using a low-cost optical particle counter and a Federal Equivalent Method Beta Attenuation Monitor.
\textit{Aerosol Sci Technol} 2020; \textbf{54}: 147--59.

\bibitem{Barkjohn2022}
Barkjohn KK, et al.
Using low-cost sensors to quantify the effects of air quality events on indoor and outdoor PM2.5 concentrations.
\textit{Sensors} 2022; \textbf{22}: 1010.

\bibitem{MacQueen1967}
MacQueen J.
Some methods for classification and analysis of multivariate observations.
In: \textit{Proc 5th Berkeley Symp on Mathematical Statistics and Probability} 1967; \textbf{1}: 281--97.

\bibitem{Rousseeuw1987}
Rousseeuw PJ.
Silhouettes: a graphical aid to the interpretation and validation of cluster analysis.
\textit{J Comput Appl Math} 1987; \textbf{20}: 53--65.

\bibitem{Gasparrini2010}
Gasparrini A, Armstrong B, Kenward MG.
Distributed lag non-linear models.
\textit{Stat Med} 2010; \textbf{29}: 2224--34.

\bibitem{Gasparrini2011}
Gasparrini A.
Distributed lag linear and non-linear models in R: the package dlnm.
\textit{J Stat Softw} 2011; \textbf{43}: 1--20.

\bibitem{Chen2016}
Chen T, Guestrin C.
XGBoost: a scalable tree boosting system.
In: \textit{Proc 22nd ACM SIGKDD Int Conf Knowledge Discovery and Data Mining} 2016: 785--94.

\bibitem{Goodfellow2016}
Goodfellow I, Bengio Y, Courville A.
\textit{Deep Learning}.
MIT Press, 2016.

\bibitem{Huang2006}
Huang G-B, Zhu Q-Y, Siew C-K.
Extreme learning machine: theory and applications.
\textit{Neurocomputing} 2006; \textbf{70}: 489--501.

\bibitem{Lundberg2017}
Lundberg SM, Lee S-I.
A unified approach to interpreting model predictions.
In: \textit{Advances in Neural Information Processing Systems} 2017; \textbf{30}: 4765--74.

\bibitem{VonElm2007}
von Elm E, Altman DG, Egger M, et al.
The Strengthening the Reporting of Observational Studies in Epidemiology (STROBE) statement: guidelines for reporting observational studies.
\textit{Lancet} 2007; \textbf{370}: 1453--57.

\bibitem{Yu2025}
Yu T, Jiang Y, Chen R, et al.
National and provincial burden of disease attributable to fine particulate matter air pollution in China, 1990--2021: an analysis of data from the Global Burden of Disease Study 2021.
\textit{Lancet Planet Health} 2025; \textbf{9}: e174--85.

\bibitem{Gopikrishnan2025}
Gopikrishnan GS, Kuttippurath J.
Impact of the National Clean Air Programme (NCAP) on the particulate matter pollution and associated reduction in human mortalities in Indian cities.
\textit{Sci Total Environ} 2025; \textbf{968}: 178787.

\bibitem{Hertzog2024}
Hertzog L, Morgan GG, Yuen C, et al.
Mortality burden attributable to exceptional PM2.5 air pollution events in Australian cities: a health impact assessment.
\textit{Heliyon} 2024; \textbf{10}: e24532.

\bibitem{Nashwan2024}
Nashwan AJ, Aldosari N, Hendy A.
Hajj 2024 heatwave: addressing health risks and safety.
\textit{Lancet} 2024; \textbf{404}: 427--28.

\end{thebibliography}

\end{document}
